% Part: normal-modal-logic
% Chapter: tableaux
% Section: completeness

\documentclass[../../../include/open-logic-section]{subfiles}

\begin{document}

\olfileid{nml}{tab}{cpl}

\olsection{تمامیت برای~\Log{K}}

\begin{explain}
  برای نشان‌دادن کامل‌بودن روش !!{tableau}ها باید نشان دهیم هرگاه هیچ
  !!{tableau} بسته‌ای برای اثبات~$\Gamma \Proves !A$ وجود نداشته
  باشد، آنگاه~$\Gamma \Entails/ !A$؛ یعنی پادمدلی وجود دارد. اما «هیچ
  !!{tableau} بسته‌ای وجود ندارد» یعنی هر راهی که برای ساختنش
  می‌آزماییم باید در بسته‌شدن شکست بخورد. نکته این است که اگر همهٔ این
  راه‌ها شکست بخورند، یک راه مشخص، \emph{نظام‌مند و فراگیر} نیز در
  بسته‌شدن شکست می‌خورد. و اگر !!{tableau} بسته‌ای وجود می‌داشت، این
  روش نظام‌مند و فراگیر بسته می‌شد. همین یک تابلو در میان شاخه‌های
  بازش همهٔ اطلاعات لازم برای تعریف پادمدل را در بر خواهد داشت. پادمدل
  حاصل از یک شاخهٔ باز در این تابلو همهٔ پیشوندهای به‌کاررفته روی آن
  شاخه را به‌عنوان جهان‌ها خواهد داشت؛ و !!a{propositional variable}~$p$
  در~$\sigma$ صادق است اگر و تنها اگر~$\sFmla{\True}{p}[\sigma]$ روی
  شاخه ظاهر شود.
\end{explain}

\begin{defn}
  شاخه‌ای در !!a{tableau} را کامل می‌نامیم اگر هرگاه شامل
  !!{formula} پیشونددار~$\sFmla{S}{!A}[\sigma]$ای باشد که قاعده‌ای بر
  آن اعمال‌پذیر است، موارد زیر را نیز شامل شود:
  \begin{enumerate}
    \item در مورد قواعد گزاره‌ایِ انباشتی، !!{formula}های پیشوندداری
      که نتایج متناظر قاعده‌اند؛
    \item در مورد قواعد گزاره‌ایِ انشعابی، یکی از !!{formula}های
      نتیجهٔ متناظر؛
    \item در مورد قواعد موجهاتیِ نیازمند پیشوند تازه، دست‌کم یک نتیجهٔ
      ممکن؛
    \item در مورد قواعد موجهاتیِ نیازمند پیشوند استفاده‌شده، نتیجهٔ
      متناظر برای هر پیشوند ظاهرشونده روی شاخه.
  \end{enumerate}
\end{defn}

\begin{explain}
برای نمونه، هر شاخهٔ کامل هرگاه~$\sFmla{\True}{!B \land
  !C}[\sigma]$ را شامل شود، $\sFmla{\True}{!B}[\sigma]$ و
$\sFmla{\True}{!C}[\sigma]$ را نیز شامل می‌شود. اگر
$\sFmla{\True}{!B \lor !C}[\sigma]$ را شامل شود، دست‌کم یکی از
$\sFmla{\True}{!B}[\sigma]$ و~$\sFmla{\True}{!C}[\sigma]$ را شامل
می‌شود. اگر \iftag{prvBox}{$\sFmla{\False}{\Box !B}[\sigma]$}
{$\sFmla{\True}{\Diamond !B}[\sigma]$} را شامل شود، برای دست‌کم یک
$n$، \iftag{prvBox}{$\sFmla{\False}{!B}[\sigma.n]$}
{$\sFmla{\True}{!B}[\sigma.n]$} را نیز شامل می‌شود. و هرگاه
\iftag{prvBox}{$\sFmla{\True}{\Box !B}[\sigma]$}
{$\sFmla{\False}{\Diamond !B}[\sigma]$} را شامل شود، برای هر~$n$ که
$\sigma.n$ روی شاخه استفاده‌شده است، \iftag{prvBox}
{$\sFmla{\True}{!B}[\sigma.n]$}
{$\sFmla{\False}{!B}[\sigma.n]$} را نیز شامل می‌شود.
\end{explain}

\begin{prop}\ollabel{prop:complete-tableau}
  هر~$\Gamma$ متناهی یک !!a{tableau} دارد که هر شاخهٔ آن کامل است.
\end{prop}

\begin{proof}
  شاخه‌ای باز در !!a{tableau}ای برای~$\Gamma$ را در نظر بگیرید. شمار
  !!{formula}های پیشونددار روی شاخه که قاعده‌ای بر آن‌ها اعمال‌پذیر
  است متناهی است. با ترتیبی ثابت (مثلاً از بالا به پایین)، برای هریک از
  این !!{formula}های پیشونددار که شرایط (1)--(4) درباره‌اش ازپیش برقرار
  نیست، قواعد اعمال‌پذیر را برای گسترش شاخه اعمال کنید. در برخی موارد
  این کار به انشعاب می‌انجامد؛ برای همهٔ !!{formula}های پیشونددار
  باقی‌مانده، قاعده را در انتهای هریک از شاخه‌های حاصل اعمال کنید. چون
  شمار !!{formula}های پیشونددار و شمار پیشوندهای استفاده‌شده روی شاخه
  متناهی است، این روش سرانجام شاخه‌هایی (احتمالاً بسیار) به دست می‌دهد
  که شاخهٔ اصلی را گسترش می‌دهند. روش را بر هریک اعمال و تکرار کنید.
  بنا بر ساخت، هر شاخه کامل است.
\end{proof}

\begin{thm}[تمامیت]
  \ollabel{thm:tableau-completeness}
  اگر~$\Gamma$ متناهی باشد و هیچ !!{tableau} بسته‌ای نداشته باشد،
  $\Gamma$ ارضاپذیر است.
\end{thm}

\begin{proof}
بنا بر گزاره، $\Gamma$ یک !!a{tableau} دارد که هر شاخهٔ آن کامل است.
چون هیچ !!{tableau} بسته‌ای ندارد، !!a{tableau}ای دارد که دست‌کم یک
شاخهٔ آن باز و کامل است. بگذارید~$\Delta$ مجموعهٔ !!{formula}های
پیشونددار روی آن شاخه و~$P(\Delta)$ مجموعهٔ پیشوندهای ظاهرشونده در آن
باشد.

مدل~$\mModel{M(\Delta)} = \tuple{P(\Delta), R, V}$ را تعریف می‌کنیم؛
جهان‌های آن پیشوندهای ظاهرشونده در~$\Delta$ هستند و رابطهٔ
دسترسی‌پذیری چنین داده می‌شود:
\[
R\sigma\sigma' \quad \text{اگر و تنها اگر} \quad
\sigma'=\sigma.n \quad \text{برای یک~$n$}
\]
و
\[
V(p) = \Setabs{\sigma}{\sFmla{\True}{p}[\sigma] \in \Delta}.
\]
با استقرا بر~$!A$ نشان می‌دهیم اگر~$\sFmla{\True}{!A}[\sigma] \in
\Delta$، آنگاه~$\mSat{M(\Delta)}{!A}[\sigma]$؛ و اگر
$\sFmla{\False}{!A}[\sigma] \in \Delta$، آنگاه
$\mSat/{M(\Delta)}{!A}[\sigma]$.

\begin{enumerate}
  \item \indcase{!A}{p}{اگر~$\sFmla{\True}{\indfrm}[\sigma] \in
    \Delta$، آنگاه~$\sigma \in V(p)$ (بنا به تعریف~$V$)، و بنابراین
    $\mSat{M(\Delta)}{\indfrm}[\sigma]$.

    اگر~$\sFmla{\False}{\indfrm}[\sigma] \in \Delta$، آنگاه
    $\sFmla{\True}{\indfrm}[\sigma] \notin \Delta$، زیرا وگرنه شاخه
    بسته می‌بود. پس~$\sigma \notin V(p)$ و بنابراین
    $\mSat/{M(\Delta)}{\indfrm}[\sigma]$.}
  \item \indcase{!A}{\lnot !B}{\iftag{probNot}{تمرین.}{اگر
      $\sFmla{\True}{\indfrm}[\sigma] \in \Delta$، آنگاه چون شاخه
      کامل است، $\sFmla{\False}{!B}[\sigma] \in \Delta$. بنا بر فرض
      استقرا، $\mSat/{M(\Delta)}{!B}[\sigma]$؛ پس
      $\mSat{M(\Delta)}{\indfrm}[\sigma]$.

      اگر~$\sFmla{\False}{\indfrm}[\sigma] \in \Delta$، آنگاه چون
      شاخه کامل است، $\sFmla{\True}{!B}[\sigma] \in \Delta$. بنا بر
      فرض استقرا، $\mSat{M(\Delta)}{!B}[\sigma]$؛ پس
      $\mSat/{M(\Delta)}{\indfrm}[\sigma]$.}}
  \item \indcase{!A}{!B \land !C}{\iftag{probAnd}{تمرین.}{اگر
      $\sFmla{\True}{\indfrm}[\sigma] \in \Delta$، آنگاه چون شاخه
      کامل است، هم~$\sFmla{\True}{!B}[\sigma] \in \Delta$ و هم
      $\sFmla{\True}{!C}[\sigma] \in \Delta$. بنا بر فرض استقرا،
      $\mSat{M(\Delta)}{!B}[\sigma]$ و
      $\mSat{M(\Delta)}{!C}[\sigma]$. پس
      $\mSat{M(\Delta)}{\indfrm}[\sigma]$.

      اگر~$\sFmla{\False}{\indfrm}[\sigma] \in \Delta$، آنگاه چون
      شاخه کامل است، یا~$\sFmla{\False}{!B}[\sigma] \in \Delta$ یا
      $\sFmla{\False}{!C}[\sigma] \in \Delta$. بنا بر فرض استقرا،
      یا~$\mSat/{M(\Delta)}{!B}[\sigma]$ یا
      $\mSat/{M(\Delta)}{!C}[\sigma]$. پس
      $\mSat/{M(\Delta)}{\indfrm}[\sigma]$.}}
  \item \indcase{!A}{!B \lor !C}{\iftag{probOr}{تمرین.}{اگر
      $\sFmla{\True}{\indfrm}[\sigma] \in \Delta$، آنگاه چون شاخه
      کامل است، یا~$\sFmla{\True}{!B}[\sigma] \in \Delta$ یا
      $\sFmla{\True}{!C}[\sigma] \in \Delta$. بنا بر فرض استقرا،
      یا~$\mSat{M(\Delta)}{!B}[\sigma]$ یا
      $\mSat{M(\Delta)}{!C}[\sigma]$. پس
      $\mSat{M(\Delta)}{\indfrm}[\sigma]$.

      اگر~$\sFmla{\False}{\indfrm}[\sigma] \in \Delta$، آنگاه چون
      شاخه کامل است، هم~$\sFmla{\False}{!B}[\sigma] \in \Delta$ و هم
      $\sFmla{\False}{!C}[\sigma] \in \Delta$. بنا بر فرض استقرا،
      هم~$\mSat/{M(\Delta)}{!B}[\sigma]$ و هم
      $\mSat/{M(\Delta)}{!C}[\sigma]$. پس
      $\mSat/{M(\Delta)}{\indfrm}[\sigma]$.}}
  \item \indcase{!A}{!B \lif !C}{\iftag{probIf}{تمرین.}{اگر
      $\sFmla{\True}{\indfrm}[\sigma] \in \Delta$، آنگاه چون شاخه
      کامل است، یا~$\sFmla{\False}{!B}[\sigma] \in \Delta$ یا
      $\sFmla{\True}{!C}[\sigma] \in \Delta$. بنا بر فرض استقرا،
      یا~$\mSat/{M(\Delta)}{!B}[\sigma]$ یا
      $\mSat{M(\Delta)}{!C}[\sigma]$. پس
      $\mSat{M(\Delta)}{\indfrm}[\sigma]$.

      اگر~$\sFmla{\False}{\indfrm}[\sigma] \in \Delta$، آنگاه چون
      شاخه کامل است، هم~$\sFmla{\True}{!B}[\sigma] \in \Delta$ و هم
      $\sFmla{\False}{!C}[\sigma] \in \Delta$. بنا بر فرض استقرا،
      هم~$\mSat{M(\Delta)}{!B}[\sigma]$ و هم
      $\mSat/{M(\Delta)}{!C}[\sigma]$. پس
      $\mSat/{M(\Delta)}{\indfrm}[\sigma]$.}}
  \item \indcase{!A}{\Box !B}{\iftag{probBox}{تمرین.}{اگر
      $\sFmla{\True}{\indfrm}[\sigma] \in \Delta$، آنگاه چون شاخه
      کامل است، برای هر~$\sigma.n$ استفاده‌شده روی شاخه، یعنی برای هر
      $\sigma' \in P(\Delta)$ که~$R\sigma\sigma'$، داریم
      $\sFmla{\True}{!B}[\sigma.n] \in \Delta$. بنا بر فرض استقرا،
      برای هر~$\sigma'$ با~$R\sigma\sigma'$،
      $\mSat{M(\Delta)}{!B}[\sigma']$. بنابراین
      $\mSat{M(\Delta)}{\indfrm}[\sigma]$.

      اگر~$\sFmla{\False}{\indfrm}[\sigma] \in \Delta$، آنگاه چون
      شاخه کامل است، برای یک~$\sigma.n$،
      $\sFmla{\False}{!B}[\sigma.n] \in \Delta$. بنا بر فرض استقرا،
      $\mSat/{M(\Delta)}{!B}[\sigma.n]$. چون~$R\sigma(\sigma.n)$،
      $\sigma'$ای وجود دارد چنان‌که
      $\mSat/{M(\Delta)}{!B}[\sigma']$. پس
      $\mSat/{M(\Delta)}{\indfrm}[\sigma]$.}}
  \item \indcase{!A}{\Diamond !B}{\iftag{probDiamond}{تمرین.}{اگر
      $\sFmla{\True}{\indfrm}[\sigma] \in \Delta$، آنگاه چون شاخه
      کامل است، برای یک~$\sigma.n$،
      $\sFmla{\True}{!B}[\sigma.n] \in \Delta$. بنا بر فرض استقرا،
      $\mSat{M(\Delta)}{!B}[\sigma.n]$. چون~$R\sigma(\sigma.n)$،
      $\sigma'$ای وجود دارد چنان‌که~$\mSat{M(\Delta)}{!B}[\sigma']$.
      پس~$\mSat{M(\Delta)}{\indfrm}[\sigma]$.

      اگر~$\sFmla{\False}{\indfrm}[\sigma] \in \Delta$، آنگاه چون
      شاخه کامل است، برای هر~$\sigma.n$ استفاده‌شده روی شاخه، یعنی
      برای هر~$\sigma' \in P(\Delta)$ که~$R\sigma\sigma'$، داریم
      $\sFmla{\False}{!B}[\sigma.n] \in \Delta$. بنا بر فرض استقرا،
      برای هر~$\sigma'$ با~$R\sigma\sigma'$،
      $\mSat/{M(\Delta)}{!B}[\sigma']$. بنابراین
      $\mSat/{M(\Delta)}{\indfrm}[\sigma]$.}}
\end{enumerate}
تابع همانی روی~$P(\Delta)$ تفسیری از پیشوندهای~$P(\Delta)$ در
$\mModel{M(\Delta)}$ است، و~$\mModel{M(\Delta)}$ نسبت به این تفسیر،
$\Gamma$ را ارضا می‌کند.
\end{proof}

\begin{prob}
اثبات \olref[nml][tab][cpl]{thm:tableau-completeness} را کامل کنید.
\end{prob}

\begin{cor}
\ollabel{cor:entailment-completeness}
اگر~$\Gamma$ متناهی باشد و~$\Gamma \Entails !A$، آنگاه
$\Gamma \Proves !A$.
\end{cor}

\begin{cor}
\ollabel{cor:weak-completeness}
اگر~$!A$ در همهٔ مدل‌ها صادق باشد، آنگاه~$\Proves !A$.
\end{cor}

\end{document}
